\documentclass[aspectratio=169,11pt,t]{beamer}
\input{../../../theme/ksa-cs-v2.tex}
\usepackage{fontawesome5}
\newcommand{\paperlink}[1]{\href{#1}{\textcolor{ksablue}{\faLink}}}
% GLM 차시 (2026-11-27): 앤드류 응 CS229 강의 노트 Part III (Generalized Linear Models) 을 덱 스타일로.
% 숫자는 data.py (assert), 그림은 make_figs.py.
\setbeamerfont{author}{size=\normalsize}
\setbeamerfont{institute}{size=\small}
\graphicspath{{figs/}{../../../assets/}}
\usetikzlibrary{positioning,arrows.meta,calc}

\renewcommand{\ksacredit}{Machine Learning and Deep Learning \,\textperiodcentered\, 포아송과 GLM}
\title{포아송과 GLM}
\subtitle{포아송 분포 · 포아송 회귀 · 일반화 선형 모델}
\author{Machine Learning and Deep Learning}
\institute{KSA 2027-Spring}
\date{Sumin Han (suminhan@ksa.hs.kr)}
\begin{document}
{\setbeamertemplate{footline}{}\begin{frame}[c,noframenumbering]\titlepage\end{frame}}

\section{질문}
% --- 오늘의 질문: 개수일 때
\begin{frame}{오늘의 질문: $y$ 가 개수일 때는?}
\vfill
\centering
{\normalsize 지금까지: \;실수 $\to$ 선형회귀 \;·\; 0 / 1 $\to$ 로지스틱회귀 \;·\; 여러 클래스 $\to$ 소프트맥스}\\[0.2em]
\begin{columns}[c]
\begin{column}{0.50\textwidth}\centering
\includegraphics[width=0.78\linewidth]{glm_v1_map.png}\\
{\scriptsize 실제 낙하 지점 887발 (Shaw 2025) · 주황 칸 0.5 km}
\end{column}
\begin{column}{0.48\textwidth}\centering
{\fontsize{40}{46}\selectfont\textcolor{ksared}{\faPlane}}\\[0.2em]
{\normalsize V-1 비행 폭탄 (Fi 103, 별명 `두들버그'): 조종사 없이 프랑스 발사대에서 날아옴}\\[0.6em]
{\LARGE 1944년 런던 시내, V-1 887발}\\[0.3em]
{\LARGE 0.5 km 칸 하나에 \kq{몇 발}?}\\[0.4em]
{\large 독일은 \kb{노리고} 쐈을까, \kb{제멋대로}였을까?}
\end{column}
\end{columns}\vskip0.3em
{\normalsize 그다음: 개수를 $x$ 로 맞히려면 직선 $z$ 에 무엇을 씌울까? \;·\; 모델들을 \kb{하나로 묶는 식}은?}
\vfill
\note{지금까지 세 가지 모델을 배웠습니다. 점수처럼 아무 실수인 y는 선형회귀, 체스 승패처럼 0 아니면 1인 y는 로지스틱회귀, 붓꽃처럼 여러 클래스 중 하나인 y는 소프트맥스 회귀로 맞혔습니다. 오늘은 개수입니다. 1944년 독일은 런던에 V-1이라는 비행 폭탄을 쏘았습니다. 비행기가 떨어뜨린 폭탄이 아니라, 폭탄 자체가 작은 제트 비행기입니다. 정식 이름은 피젤러 Fi 103, V는 보복 무기라는 독일어의 머리글자이고, 영국 사람들은 붕붕거리는 소리 때문에 두들버그, 버즈 밤이라고 불렀습니다. 조종사 없이 프랑스 해안의 발사대에서 런던 쪽으로 쏘면, 연료가 떨어지는 곳에서 엔진이 꺼지고 떨어졌습니다. 런던 시의회 구역에만 887발이 떨어졌습니다. 왼쪽 지도의 빨간 점이 실제로 폭탄이 떨어진 887곳입니다. 전쟁 뒤 런던 시의회가 만든 폭격 피해 지도에서 옮긴 좌표입니다. 주황 네모 안을 0.5킬로미터 칸으로 나눠 칸마다 몇 발인지 셉니다. 0.5킬로미터 칸 하나에는 몇 발이 떨어졌을까요? 0발, 1발, 2발로 세는 값, 개수입니다. 그리고 당시 사람들은 궁금했습니다. 독일이 특정 구역을 노리고 쏜 걸까, 아니면 제멋대로 떨어진 걸까? 오늘 이 질문에 답해 봅니다. 그다음에는 개수를 x로 맞히는 모델, 직선 z에 무엇을 씌우고 손실은 무엇으로 할지, 그리고 마지막에 이 모델들을 하나로 묶는 식이 있는지 봅니다.}
\end{frame}

\section{포아송 분포}
% --- 개수 예시 여럿
\begin{frame}{개수 데이터는 비슷한 모양으로 흩어진다}
\vfill
\centering\includegraphics[width=0.84\linewidth]{glm_pois_examples.png}\par\vskip0.2em
{\large 막대 $=$ 실제 비율 \;·\; 검은 점 $=$ 같은 평균 $\lambda$ 의 \kb{포아송} \;$\to$\; 셋 다 거의 겹친다}\\[0.3em]
{\large 공통점: 드물게 일어나는 일을 \kb{정해진 칸} (한 구역 · 부대 1년 · 운전자 1명의 1년) 안에서 센다}\\[0.3em]
{\small 폭탄 = Clarke (1946) 런던 남부 576구역 \quad 기병대 = Bortkiewicz (1898) \quad 보험 = 프랑스 freMTPL2, 169,349계약}
\vfill
\note{개수를 세는 데이터를 세 가지 봅시다. 왼쪽이 오늘의 질문, 런던 폭탄입니다. 런던 남부를 0.25제곱킬로미터짜리 576개 구역으로 나누고 구역마다 떨어진 V-1 폭탄 수를 셌습니다. 229구역은 0발, 211구역은 1발, 93구역은 2발이었고, 평균은 구역당 0.93발입니다. 가운데는 포아송 분포의 가장 유명한 예입니다. 1898년 보르트키에비치가 프로이센 기병 10개 부대의 20년 기록, 200 부대-년에서 한 해에 말에 차여 숨진 병사 수를 셌습니다. 평균은 0.61명입니다. 오른쪽은 뒤에서 모델을 만들 자동차 보험 데이터입니다. 프랑스 자동차 보험 계약 169,349개에서 운전자가 1년에 사고를 몇 번 청구했는지 셌고, 평균 0.070번입니다. 막대는 실제 비율, 검은 점은 같은 평균을 가진 포아송 분포인데 셋 다 거의 겹칩니다. 공통점은 드물게 일어나는 일을 한 구역, 부대의 1년, 운전자 한 명의 1년처럼 정해진 칸 안에서 센다는 것입니다. 이렇게 셌을 때 나오는 분포가 포아송입니다.}
\end{frame}

% --- 포아송의 발명
\begin{frame}{포아송은 어디서 왔나: 드문 일을 많이 시도하면}
\vfill
\centering\includegraphics[width=0.74\linewidth]{glm_pois_origin.png}\par\vskip0em
{\large 포커 5장: 판 수 $\times$ 확률 $=$ 평균 1.27번으로 같게 \;$\to$\; \kb{판은 많고 패는 드물수록} 포아송}\\[0.3em]
{\large\[ \underbrace{\binom{n}{y} p^{y} (1-p)^{n-y}}_{n \text{ 판 중 } y \text{ 번 (이항)}} \;\xrightarrow[\;np \,=\, \lambda\;]{\;n \to \infty,\; p \to 0\;}\; \frac{\lambda^{y} e^{-\lambda}}{y!} \]}\vskip-0.3em
{\small 포아송 (Poisson, 1837) 이 재판 판결 확률을 연구하다 얻은 식 \quad 보르트키에비치 (1898): \kb{작은 수의 법칙}}
\vfill
\note{포아송 분포는 어디서 왔을까요? 포커로 봅시다. 카드 5장을 받을 때 원페어가 나올 확률은 42퍼센트, 트리플은 2.1퍼센트, 풀하우스는 0.144퍼센트입니다. 원페어는 3판, 트리플은 60판, 풀하우스는 880판을 하면, 세 경우 모두 평균 1.27번 나옵니다. 몇 번 나오는지는 동전 던지기를 여러 번 한 이항분포입니다. 막대가 그 실제 확률이고, 검은 점은 같은 평균 1.27의 포아송입니다. 원페어 3판은 최대 3번이라 포아송과 꽤 다릅니다. 트리플 60판은 거의 겹치고, 풀하우스 880판은 완전히 겹칩니다. 판은 많고 한 판에 나올 확률은 작을수록 이항분포가 포아송이 됩니다. 식으로는 평균 n 곱하기 p를 람다로 고정하고 n을 끝없이 키운 극한입니다. 프랑스 수학자 포아송이 1837년에 재판에서 잘못된 판결이 나올 확률을 연구하는 책에서 이 식을 얻었고, 1898년 보르트키에비치가 기병대 자료로 드문 일의 개수가 이 식을 따른다는 것을 보이고 작은 수의 법칙이라고 불렀습니다.}
\end{frame}

% --- 연결고리
\begin{frame}{그래서 포아송: 기회는 많고, 한 번 한 번은 드물다}
\vfill
\centering
\renewcommand{\arraystretch}{1.3}
{\large\begin{tabular}{l|l|l|c}
센 것 & 기회 (동전 $n$ 번) & 한 번의 확률 $p$ & 평균 $\lambda = np$\\ \hline
포커: 풀하우스 & 880판 & 풀하우스: 0.144\% & 1.27번\\
런던: 구역의 폭탄 & 폭탄 537발 & 이 구역에: $1/576$ & 0.93발\\
기병대: 숨진 병사 & 병사 수백 명 & 말에 차여 숨짐 (드묾) & 0.61명\\
{\faCar}\; 보험: 1년 사고 & 1년 365일 & 그날 사고 (드묾) & 0.070번\\
{\faPhone}\; 콜센터: 1분 전화 & 고객 수천 명 & 이 1분에 전화함 & $-$\\
{\faCookie}\; 쿠키: 초코칩 & 초코칩 수백 개 & 이 쿠키에 들어감 & $-$
\end{tabular}}\\[0.8em]
{\large $n$ 은 크고 \;·\; $p$ 는 작고 \;·\; 한 번 한 번이 \kb{서로 상관없으면} \;$\to$\; 개수는 \kb{포아송}}\\[0.3em]
{\small 확인: 537발 이항으로 기대한 구역 수 227 · 212 · 99 · 31 · 7 \quad 실제 229 · 211 · 93 · 35 · 7 (0$\sim$4발)}
\vfill
\note{이제 앞에서 본 예시들이 왜 포아송인지 연결됩니다. 모두 포커의 풀하우스처럼, 기회는 아주 많은데 한 번 한 번은 드문 동전 던지기입니다. 런던 폭탄은 더 분명합니다. 537발이 떨어졌고, 한 발이 특정 구역에 떨어질 확률은 576분의 1입니다. 동전 537번을 던지는 이항분포로 계산하면 폭탄이 0발인 구역 227개, 1발 212개, 2발 99개, 3발 31개, 4발 7개를 기대하는데, 실제는 229, 211, 93, 35, 7개였습니다. 기병대는 부대마다 병사가 수백 명이고, 한 사람이 한 해에 말에 차여 숨질 확률은 아주 작습니다. 보험도 같습니다. 운전자 한 명은 1년 365일 매일 동전을 던지는데, 그날 사고가 날 확률은 아주 작습니다. 그래서 1년 사고 횟수가 평균 0.070번인 포아송입니다. 콜센터도 고객 수천 명이 각자 이 1분에 전화할 확률은 작고, 서로 짜고 걸지 않으니 1분 동안 걸려 오는 전화 수가 포아송입니다. 초코칩 쿠키도 반죽을 잘 섞으면 초코칩 하나하나가 이 쿠키에 들어갈 확률이 작아서, 쿠키 하나의 초코칩 수가 포아송입니다. 조건은 세 가지입니다. 기회가 많고, 한 번의 확률이 작고, 한 번 한 번이 서로 상관없어야 합니다. 마지막 조건이 깨지면, 예를 들어 비 오는 날 따릉이를 다 같이 안 타는 것처럼 함께 움직이면, 포아송보다 넓게 퍼집니다.}
\end{frame}

% --- 흩뿌린 점
\begin{frame}{서로 상관없어야: 점을 흩뿌리고 칸마다 세면}
\vfill
\centering\includegraphics[width=0.64\linewidth]{glm_pois_scatter.png}\par\vskip0em
{\large 칸마다 평균 2개, \kb{분산} 1.9 \;·\; 0.0 \;·\; 10.9 \;$\to$\; 포아송은 \kb{분산 $=$ 평균}}\\[0.2em]
{\large 런던 폭탄 $\approx$ \kb{제멋대로} $\to$ 노린 게 아니었다}\\[0.1em]
{\normalsize 포아송이 아닌 것: 전염병 확진자 (서로 옮긴다) \;·\; 출근길 지하철 승객 (한때 몰린다)}
\vfill
\note{앞 쪽의 세 번째 조건, 서로 상관없어야 한다는 것을 그림으로 봅시다. 같은 넓이에 점 192개를 세 가지 방식으로 뿌리고, 96개 칸마다 점이 몇 개인지 셉니다. 평균은 셋 다 칸당 2개입니다. 왼쪽은 점마다 제멋대로, 서로 상관없이 떨어졌습니다. 칸마다 센 개수가 포아송과 거의 겹치고, 분산도 1.9로 평균 2와 거의 같습니다. 가운데는 점들이 서로 밀어내서 고르게 퍼졌습니다. 모든 칸에 딱 2개씩, 분산이 0이고 포아송보다 좁습니다. 오른쪽은 점들이 서로 끌어당겨 뭉쳤습니다. 빈 칸이 아주 많고 어떤 칸은 8개가 넘습니다. 분산이 10.9로 평균보다 훨씬 크고, 포아송보다 넓습니다. 포아송은 분산이 평균과 같은 분포입니다. 이제 처음 질문에 답할 수 있습니다. 독일이 특정 구역을 노렸다면 폭탄은 오른쪽처럼 뭉쳤어야 합니다. 그런데 런던 폭탄은 포아송과 딱 맞았습니다. 왼쪽처럼 제멋대로 떨어졌다는 뜻이고, 영국은 이 분석으로 V-1이 특정 구역을 정밀하게 노릴 수 없다고 판단했습니다. 반대로 전염병 확진자 수는 한 사람이 주변에 옮기니 서로 상관있어서 뭉치고, 출근길 지하철 승객은 특정 시간에 몰려서 포아송보다 넓게 퍼집니다.}
\end{frame}

% --- 실제 지도 시뮬레이션
\begin{frame}{실제 지도로 확인: 제멋대로 1만 번 떨어뜨려 보면}
\vfill
\centering\includegraphics[width=0.86\linewidth]{glm_v1_sim.png}\par\vskip0.2em
{\large 칸 432개에 383발: 분산 $\div$ 평균 $=$ \kb{1.10} \;(제멋대로면 0.88$\sim$1.14)}\\[0.2em]
{\large 칸 수도 모두 범위 안 \;$\to$\; 제멋대로와 \kb{구분이 안 된다} $=$ 노리지 못했다}
\vfill
\note{정말 포아송인지 실제 낙하 지점으로 확인해 봅시다. 지도의 주황 네모, 템스강 남쪽 가로 12킬로미터 세로 9킬로미터를 클라크처럼 0.5킬로미터 칸, 432칸으로 나누면 폭탄이 383발입니다. 빨간 막대가 칸별 실제 개수입니다. 0발인 칸 190개, 1발 141개, 2발 68개, 3발 26개, 4발 7개입니다. 검은 점은 같은 평균 0.89의 포아송입니다. 이제 컴퓨터로 383발을 432칸에 제멋대로, 칸마다 같은 확률로 떨어뜨리기를 1만 번 해 봅니다. 회색 막대가 그 1만 번의 95퍼센트 범위입니다. 실제 칸 수는 모두 이 범위 안에 들어옵니다. 오른쪽은 1만 번 각각의 분산 나누기 평균입니다. 제멋대로면 1 근처, 0.88에서 1.14 사이에 들어오고, 실제는 1.10으로 그 안에 있습니다. 실제 폭탄은 제멋대로 떨어뜨린 것과 구분이 안 됩니다. 독일이 특정 칸을 노렸다면 폭탄이 몇 칸에 뭉쳐서, 빈 칸과 많이 맞은 칸이 범위보다 훨씬 많았을 것입니다. 클라크가 1946년에 576칸으로 얻은 결론을 실제 낙하 지점으로 다시 확인한 셈입니다. 이렇게 포아송은 제멋대로라는 기준선이고, 데이터가 그 기준에서 벗어났는지를 시뮬레이션으로 잴 수 있습니다. 참고로 격자를 어디에 놓느냐에 따라 숫자는 조금씩 달라집니다.}
\end{frame}

% --- 포아송 분포
\begin{frame}{개수의 분포: 포아송}
\vfill
\centering
{\Large\[ P(y) = \frac{\lambda^{y}\,e^{-\lambda}}{y!} \qquad y = 0, 1, 2, \dots \qquad \text{평균} = \lambda \]}\vskip-0.2em
\includegraphics[width=0.5\linewidth]{glm_pois_pmf.png}\par\vskip0em
{\large 평균 $\lambda$ 하나가 모양을 정한다 \;·\; $\lambda$ 가 클수록 \kb{오른쪽으로, 넓게}}
\vfill
\note{개수에 쓰는 분포가 포아송입니다. 평균이 람다일 때 y번 나올 확률은 람다의 y제곱 곱하기 e의 마이너스 람다 나누기 y 팩토리얼입니다. 평균 람다 하나만 정하면 모양이 다 정해집니다. 그림은 뒤에서 학습할 보험 모델이 예측한 세 경우입니다. 초록은 보너스-말러스 등급 50, 오래 사고 없이 탄 운전자로 평균 0.06번, 1년에 사고가 없을 확률이 94퍼센트입니다. 회색은 처음 가입한 등급 100, 평균 0.22번입니다. 빨강은 사고 이력이 많은 등급 150, 평균 0.79번으로 1번, 2번도 흔합니다. 람다가 커질수록 오른쪽으로 가면서 넓게 퍼집니다. 그러니 x로 람다만 잘 맞히면 됩니다.}
\end{frame}

\section{포아송 회귀}
% --- 개수에 직선
\begin{frame}{개수를 $x$ 로 맞히기: 보험 사고 횟수에 직선을 그으면}
\vfill
\centering\includegraphics[width=0.53\linewidth]{glm_claims_line.png}\par\vskip0.1em
{\large 직선: 청구 $= -0.097 + 0.00306\,x$ \quad 등급 140$\sim$149: 실제 \kb{0.72번}, 직선은 0.35번}\\[0.2em]
{\large 실제 평균은 등급이 오를수록 \kb{곱해서} 늘어난다 (휜다) \;$\to$\; 직선은 못 따라간다}\\[0.2em]
{\small $x$ = 보너스-말러스 등급: 처음 100, 무사고 1년마다 $\times 0.95$ (최저 50), 사고 1번마다 $\times 1.25$}
\vfill
\note{이제 개수를 x로 맞혀 봅시다. 폭탄 데이터에는 구역마다 센 개수만 있고 x가 없으니, x가 있는 자동차 보험 데이터를 씁니다. 보험 회사는 운전자가 내년에 사고를 몇 번 낼지 맞혀야 보험료를 정합니다. x는 보너스-말러스 등급입니다. 우리나라 자동차 보험의 할인 할증 등급과 같은 것으로, 처음 가입하면 100이고, 사고 없이 1년을 보내면 0.95배로 내려가 최저 50까지, 사고를 한 번 내면 1.25배로 올라갑니다. 등급이 높을수록 사고 이력이 많은 운전자입니다. 검은 점은 등급 10점 구간마다 실제 평균 사고 횟수이고, 점이 클수록 계약이 많습니다. 대부분은 등급 50, 오래 무사고인 운전자입니다. 1주차처럼 직선을 그으면 청구 횟수는 마이너스 0.097 더하기 0.00306 곱하기 등급입니다. 등급이 낮은 쪽은 괜찮은데, 등급이 높은 쪽에서 크게 틀립니다. 등급 140에서 149인 50계약은 실제로 평균 0.72번 사고를 냈는데 직선은 0.35번입니다. 실제 평균은 등급이 오를수록 더하는 게 아니라 곱해서 늘어나 위로 휩니다. 그리고 직선은 등급 32 아래로 가면 음수 사고를 예측합니다. 사고를 마이너스 번 낼 수는 없죠. 로지스틱회귀 때 0과 1을 넘는 직선 대신 시그모이드를 썼듯이, 여기도 다른 모양이 필요합니다.}
\end{frame}

% --- 포아송 회귀 모델
\begin{frame}{포아송 회귀: 평균 $\lambda$ 를 직선으로}
\vfill
\centering
{\large $\lambda = w_0 + w_1 x$ 로 두면? \;$\to$\; 휘지 못하고, 끝에서는 \kq{음수 사고} (평균은 0 이상이어야)}\\[0.8em]
{\LARGE\[ \hat y = \lambda = e^{\,z}, \qquad z = w_0 + w_1 x \]}\vskip0.4em
{\large $e^{z}$ 는 \kb{언제나 양수} \;·\; $x$ 가 1 늘면 $z$ 는 $w_1$ 늘고 $\hat y$ 는 \kb{$e^{w_1}$ 배}}\\[0.6em]
{\large 로지스틱과 같은 구조: 직선 $z$ 에 함수 하나 ($\sigma$ 대신 $e^{z}$)}\\[0.3em]
{\small 거꾸로 쓰면 $z = \log\lambda$ \;(평균의 로그가 직선) \quad 로지스틱은 $z = \log\frac{p}{1-p}$ (오즈의 로그가 직선)}
\vfill
\note{포아송의 평균 람다를 x로 맞히는 모델을 만들어 봅시다. 1주차처럼 람다를 바로 직선으로 두면, 곱해서 늘어나는 모양을 못 따라가고, 끝에서는 0 아래로 내려가 음수 사고를 예측합니다. 평균 사고 횟수는 0 이상이어야 하죠. 그래서 직선 z에 e를 씌웁니다. 와이햇은 람다, 람다는 e의 z입니다. e의 z는 z가 아무리 작아도 언제나 양수입니다. 그리고 x가 1 늘면 z가 w1 늘고, 와이햇은 e의 w1 배가 됩니다. 더해 가는 게 아니라 곱해 갑니다. 로지스틱회귀와 똑같은 구조입니다. 직선 z에 함수 하나를 씌우는데, 시그모이드 대신 e의 z입니다. 거꾸로 쓰면 z는 log 람다, 평균의 로그가 직선입니다. 로지스틱회귀에서 오즈의 로그가 직선이었던 것과 같은 읽는 법입니다.}
\end{frame}

% --- 포아송 손실 · 기울기
\begin{frame}{포아송 회귀의 손실과 기울기}
\vfill
\centering
{\large\begin{align*}
\text{loss} &= -\log P(y) = -\log\frac{\hat y^{\,y} e^{-\hat y}}{y!} = \hat y - y\log\hat y + \log y! && \text{로지스틱처럼 } -\log P\\[0.3em]
\frac{\partial\,\text{loss}}{\partial z} &= \frac{\partial}{\partial z}\big(e^{z} - y\,z\big) = e^{z} - y = \textcolor{ksared}{\hat y - y} && \hat y = e^{z},\; \log\hat y = z\\[0.3em]
\frac{\partial\,\text{loss}}{\partial w_1} &= (\hat y - y)\,x && z = w_0 + w_1 x
\end{align*}}\vskip0.4em
{\large 예: 등급 100 ($\hat y = 0.22$) 인 운전자가 1년에 1번 사고 \;$\to$\; $\hat y - y = -0.78$ \;$\to$\; 예측을 \kb{올린다}}\\[0.5em]
{\large 기울기가 \kq{로지스틱회귀와 똑같은} $(\hat y - y)\,x$ \;$\to$\; 경사하강법 코드에서 \texttt{sigmoid} 를 \texttt{exp} 로만}
\vfill
\note{손실은 로지스틱회귀 때처럼 마이너스 log 확률입니다. 실제로 나온 사고 횟수 y의 확률이 클수록 손실이 작습니다. 포아송 확률에 마이너스 log를 씌우면 와이햇 빼기 y log 와이햇 더하기 log y 팩토리얼입니다. 마지막 항은 w와 상관없어서 미분하면 사라집니다. 와이햇은 e의 z, log 와이햇은 z이니, 손실은 e의 z 빼기 y z이고, z로 미분하면 e의 z 빼기 y, 곧 와이햇 빼기 y입니다. 연쇄 법칙으로 x를 곱하면 w1으로 미분한 값은 와이햇 빼기 y 곱하기 x. 예를 들어 등급 100인 운전자에게 0.22번을 예측했는데 실제로 1번 사고를 냈다면, 와이햇 빼기 y가 마이너스 0.78이라 예측을 올리는 쪽으로 w가 바뀝니다. 놀라운 점은 이 기울기가 로지스틱회귀와 똑같다는 것입니다. 그래서 로지스틱회귀 경사하강법 코드에서 시그모이드를 exp로 바꾸기만 하면 포아송 회귀가 됩니다. 이 걸음을 실제로 되풀이해 봅시다.}
\end{frame}

% --- 포아송 학습 과정
\begin{frame}{학습: 경사하강법을 되풀이하면}
\vfill
\centering
\renewcommand{\arraystretch}{1.45}
{\large\begin{tabular}{l|l}
\kb{시작} & $w_0 = \log 0.070 = -2.66,\; w_1 = 0$ \;$\to$\; 모두 1년 0.070번으로 예측\\
\kb{한 걸음} & 169,349계약마다 $\hat y = e^{\,w_0 + w_1 x}$ \;$\to$\; 오차 $\hat y - y$ \;$\to$\; 평균 낸 기울기\\
 & $w_0 \leftarrow w_0 - \alpha\cdot\text{평균}(\hat y - y), \quad w_1 \leftarrow w_1 - \alpha\cdot\text{평균}\big((\hat y - y)\,x\big)$\\
\kb{첫 걸음 값} & $\text{평균}\big((\hat y - y)\,x\big) = -0.035$ \;($x$ 는 10점 단위) \;$\to$\; $w_1$ 이 \kb{커진다}\\
 & 등급 높은 사람이 실제로 더 많이 사고 $\to$ $\hat y - y < 0$ 이 큰 $x$ 에 몰린다\\
\kb{끝} & 20000번 되풀이 \;$\to$\; $w_0 = -4.14,\; w_1 = 0.26$ (10점당) \;(기울기 $\approx 0$)
\end{tabular}}\\[0.8em]
{\large 로지스틱회귀 코드와 \kb{한 군데만} 다르다: \;\texttt{sigmoid(z)} $\to$ \texttt{exp(z)}}
\vfill
\note{그럼 포아송 회귀는 실제로 어떻게 학습될까요? 로지스틱회귀와 똑같이 경사하강법입니다. 시작은 w0를 평균 사고 횟수 0.070의 로그, 마이너스 2.66으로, w1을 0으로 둡니다. 그러면 등급과 상관없이 모든 운전자를 1년 0.070번으로 예측합니다. 한 걸음은 이렇습니다. 169,349계약 각각에서 와이햇, e의 w0 더하기 w1 x를 계산하고, 실제 사고 횟수를 빼서 오차 와이햇 빼기 y를 구합니다. w0는 오차의 평균만큼, w1은 오차 곱하기 x의 평균만큼 학습률을 곱해 빼 줍니다. 첫 걸음에서 오차 곱하기 x의 평균은 등급 10점 단위로 마이너스 0.035입니다. 사고를 낸 운전자는 오차가 음수인데, 그런 운전자가 등급이 높은 쪽, x가 큰 쪽에 몰려 있어서 곱의 평균이 음수입니다. 기울기가 음수이니 w1을 키웁니다. 등급이 높을수록 사고가 많다는 쪽으로 곡선이 기울어지는 겁니다. 이 걸음을 2만 번 되풀이하면 기울기가 거의 0이 되고, w0 마이너스 4.14, w1은 등급 10점당 0.26, 1점당 0.026에서 멈춥니다. 다음 쪽 곡선이 바로 이 값입니다. 코드로 보면 로지스틱회귀의 sigmoid를 exp로 바꾼 것뿐입니다.}
\end{frame}

% --- 포아송 회귀 결과
\begin{frame}{보험 169,349계약으로 학습하면}
\vfill
\centering\includegraphics[width=0.56\linewidth]{glm_claims_pois.png}\par\vskip0.2em
{\large 청구 $= e^{\,-4.137 + 0.0260\,x}$ \quad 등급 50: 0.06번 \;·\; 100: 0.22번 \;·\; 150: \kb{0.79번} (직선은 0.36)}\\[0.2em]
{\large 등급 10점 오를 때마다 \kb{$\times 1.30$} ($e^{0.26}$) \quad 언제나 양수}
\vfill
\note{점선은 1주차처럼 그은 직선이고, 파란 곡선이 포아송 회귀입니다. 169,349계약으로 학습하면 1년 사고 횟수는 e의 마이너스 4.137 더하기 0.0260 곱하기 등급입니다. 오래 무사고인 등급 50은 0.06번, 처음 가입한 100은 0.22번, 사고 이력이 많은 150은 0.79번을 예측합니다. 등급이 10점 오를 때마다 e의 0.26, 1.30배씩 늘어납니다. 그래서 더해 가는 직선이 아니라 곱해 가는 지수 곡선이고, 검은 점, 실제 평균을 직선보다 잘 따라갑니다. 직선은 등급 150에서 0.36번으로 너무 낮게 봤습니다. 보험 회사는 이런 예측으로 보험료를 정합니다. 학습은 로지스틱회귀 코드에서 시그모이드를 exp로 바꾼 경사하강법으로 했고, sklearn의 PoissonRegressor와 같은 답이 나옵니다.}
\end{frame}

\section{GLM}
% --- 분포 네 개
\begin{frame}{데이터마다 맞는 분포가 있다}
\vfill
\centering
\includegraphics[width=0.95\linewidth]{glm_dists}\\[0.1em]
\renewcommand{\arraystretch}{1.2}
{\large\begin{tabular}{cccc}
\makebox[0.21\linewidth]{$\frac{1}{\sqrt{2\pi}\sigma}\,e^{-(y-\mu)^2/2\sigma^2}$} &
\makebox[0.21\linewidth]{$\phi^{\,y}(1-\phi)^{1-y}$} &
\makebox[0.21\linewidth]{$\lambda^{y}e^{-\lambda}/\,y!$} &
\makebox[0.21\linewidth]{$p(y = k) = \phi_k$}\\
$\sigma = 4.4$ 점 & $\phi = 363/500 = 0.73$ & $\lambda = 0.056$ 번 & $\phi = (0,\ 0.38,\ 0.62)$
\end{tabular}}\\[0.4em]
{\small 막대 $=$ 실제 데이터, 검은 선 $=$ 같은 평균의 분포 }
\vfill
\note{네 데이터에서 y마다 어떤 모양으로 흩어지는지, 분포를 맞춰 비교해 봅시다. 막대가 실제 데이터, 검은 선이 같은 평균을 가진 분포입니다. 점수는 직선 예측에서 벗어난 정도를 30명 모아 보면 종 모양, 정규분포입니다. 표준편차는 4.4점입니다. 체스는 레이팅이 200점 앞선 근처 500판을 모으면 백이 363판 이겼습니다. 1일 확률 파이가 0.73인 베르누이입니다. 보험은 오래 무사고인 등급 50의 12만 3천 계약에서 1년 사고 횟수입니다. 평균 0.056번인 포아송과 거의 겹치고, 분산도 0.058로 평균과 거의 같습니다. 붓꽃은 꽃잎 길이 4.8에서 5.0센티미터인 13송이를 모으면 setosa 0, versicolor 5, virginica 8송이, 범주형 분포입니다. 정규분포는 선형회귀, 베르누이는 로지스틱회귀, 포아송은 포아송 회귀, 범주형은 소프트맥스 회귀가 맡습니다. 분포 하나에 모델 하나씩입니다. 그런데 이 네 모델에는 묘한 공통점이 있습니다.}
\end{frame}

% --- 분포만 바꾸면 (한 표)
\begin{frame}{분포만 바꾸면 모델이 나온다: 기울기는 모두 $\hat y - y$}
\vfill
\centering
\renewcommand{\arraystretch}{1.5}
{\large\begin{tabular}{l|l|c|c|c}
분포 ($y$) & 모델 & 예측 $\hat y$ & 손실 $= -\log p$ & 기울기 $\partial L/\partial z$\\ \hline
정규 (실수) & 선형회귀 & $z$ & $\tfrac12 (y - \hat y)^2$ & \kq{$\hat y - y$}\\
베르누이 (0 / 1) & 로지스틱회귀 & $\sigma(z)$ & $-\log\hat y_{\text{정답}}$ (이진) & \kq{$\hat y - y$}\\
포아송 (개수) & 포아송 회귀 & $e^{z}$ & $\hat y - y\log\hat y$ & \kq{$\hat y - y$}\\
범주형 (클래스) & 소프트맥스 & $\mathrm{softmax}(z)$ & $-\log\hat y_{\text{정답}}$ & \kq{$\hat y - y$}
\end{tabular}}\\[0.8em]
{\large \kb{GLM} $=$ $y$ 의 분포 고르기 + 예측 $=$ 그 분포의 평균 + 직선 $z = w_0 + w_1 x$}\\[0.3em]
{\large 직선 $z$ 는 그대로, 바깥 함수와 손실만 분포 따라 바뀐다}\\[0.3em]
{\small 기울기가 모두 같은 이유: 네 분포가 모두 \kb{지수족} (같은 꼴의 식) $\to$ 보충 덱}
\vfill
\note{지금까지를 한 표로 모읍니다. y가 실수면 정규분포를 고르고 선형회귀, 예측은 z 그대로, 손실은 오차 제곱입니다. 0과 1이면 베르누이, 로지스틱회귀, 시그모이드와 교차 엔트로피입니다. 개수면 포아송, 포아송 회귀, e의 z와 포아송 손실입니다. 여러 클래스면 범주형, 소프트맥스 회귀, 소프트맥스와 교차 엔트로피입니다. 손실은 모두 그 분포의 마이너스 log 확률입니다. 이렇게 y의 분포를 고르고, 예측을 그 분포의 평균으로, 직선 z를 그 안에 넣는 요리법을 일반화 선형 모델, GLM이라고 부릅니다. 직선, Linear는 그대로 두고 정규분포가 아닌 y까지 다루게 넓혔다는 뜻에서 일반화, Generalized입니다. 그리고 손실을 z로 미분하면 넷 다 와이햇 빼기 y입니다. 우연이 아니라, 네 분포가 모두 지수족이라는 같은 꼴의 식으로 쓰이기 때문입니다. 그 이유는 보충 덱에서 유도합니다. 수업에서는 분포를 고르면 출력 함수와 손실이 짝으로 따라 나오고, 기울기는 언제나 와이햇 빼기 y라는 것만 기억하면 됩니다.}
\end{frame}

\section{신경망}
% --- 신경망의 마지막 층
\begin{frame}{신경망으로 이어진다: 마지막 층이 GLM}
\vfill
\centering
\begin{tikzpicture}[node distance=1.1cm, every node/.style={font=\large},
  box/.style={draw, rounded corners=3pt, thick, minimum height=1.0cm, minimum width=1.6cm, align=center}]
\node[box] (x) {입력\\$x$};
\node[box, right=of x, fill=gray!10] (h) {숨은 층들\\특징 $h$};
\node[box, right=of h, draw=ksablue, fill=ksablue!8] (z) {$z = w \cdot h + b$};
\node[box, right=of z, draw=ksablue, fill=ksablue!8] (g) {$\hat y = g(z)$};
\node[box, right=of g] (L) {손실\\$-\log p$};
\draw[-{Stealth},thick] (x) -- (h); \draw[-{Stealth},thick] (h) -- (z); \draw[-{Stealth},thick] (z) -- (g); \draw[-{Stealth},thick] (g) -- (L);
\draw[ksablue, thick, dashed, rounded corners] ([xshift=-6pt,yshift=8pt]z.north west) rectangle ([xshift=6pt,yshift=-8pt]g.south east);
\node[ksablue, font=\large\bfseries, above=10pt of $(z.north)!0.5!(g.north)$] {마지막 층 $=$ GLM};
\draw[-{Stealth}, ultra thick, ksared] ([yshift=-14pt]L.south) to[bend left=18] node[below, font=\large\bfseries] {역전파의 시작: $\hat y - y$} ([yshift=-14pt]z.south);
\end{tikzpicture}\vskip1.0em
{\large 분류 신경망의 끝 $=$ 로지스틱 / 소프트맥스 회귀 \quad 회귀 신경망의 끝 $=$ 선형회귀}\\[0.3em]
{\large 앞 층들은 $x$ 대신 \kb{좋은 특징 $h$} 를 만드는 일을 한다}
\vfill
\note{GLM은 딥러닝으로 그대로 이어집니다. 신경망은 입력 x를 숨은 층들에 통과시켜 특징 h를 만들고, 마지막 층에서 z 이퀄 w 점곱 h 더하기 b를 계산한 뒤 함수 g를 씌웁니다. 이 마지막 층이 바로 오늘 본 GLM입니다. 분류하는 신경망의 끝은 로지스틱회귀나 소프트맥스 회귀이고, 숫자를 맞히는 신경망의 끝은 선형회귀입니다. 앞의 층들이 하는 일은 x 대신 GLM이 잘 쓸 수 있는 좋은 특징 h를 만드는 것입니다. 그리고 학습할 때 손실을 z로 미분한 값은, 오늘 본 대로 와이햇 빼기 y입니다. 역전파는 바로 이 와이햇 빼기 y에서 시작해서 앞 층으로 전해집니다. 7주차 역전파에서 다시 만납니다.}
\end{frame}

% --- 분포 하나에서 출력 함수와 손실이 둘 다
\begin{frame}{분포는 activation 도 loss 도 아니다: 둘 다 분포에서 나온다}
\vfill
\centering
\begin{tikzpicture}[every node/.style={font=\large},
  box/.style={draw, rounded corners=3pt, thick, minimum height=1.0cm, minimum width=3.4cm, align=center}]
\node[box, draw=ksablue, fill=ksablue!8, font=\Large\bfseries] (d) {분포 하나 고르기};
\node[box, below left=1.3cm and 0.6cm of d] (g) {출력 함수 $g(z)$\\{\normalsize 마지막 층의 activation}};
\node[box, below right=1.3cm and 0.6cm of d] (l) {손실 함수\\{\normalsize loss}};
\draw[-{Stealth}, thick] (d.south) -- node[left=6pt, font=\normalsize] {평균을 $z$ 로 풀기} (g.north);
\draw[-{Stealth}, thick] (d.south) -- node[right=6pt, font=\normalsize] {$-\log p$} (l.north);
\node[below=0.25cm of g, font=\normalsize] {예: 베르누이 $\to$ 시그모이드 $\sigma(z)$};
\node[below=0.25cm of l, font=\normalsize] {예: 베르누이 $\to$ 이진 교차 엔트로피};
\end{tikzpicture}\vskip1.2em
{\large 같은 분포에서 나왔으니 \kb{짝}으로 다닌다 \;$\to$\; 미분하면 $\hat y - y$}\\[0.3em]
{\small 숨은 층의 activation (ReLU 등) 은 분포와 상관없다}
\vfill
\note{자주 헷갈리는 질문이 있습니다. 정규분포, 베르누이 같은 분포는 activation 함수일까요, 손실 함수일까요? 둘 다 아닙니다. 분포는 y가 이렇게 나온다는 가정입니다. 그리고 이 가정 하나에서 둘이 다 나옵니다. 분포의 평균을 z로 풀면 출력 함수 g가 나오고, 이게 신경망 마지막 층의 activation입니다. 분포에 마이너스 log를 씌우면 손실 함수가 나옵니다. 베르누이를 고르면 출력은 시그모이드, 손실은 이진 교차 엔트로피입니다. 같은 분포에서 나왔으니 둘은 짝으로 다니고, 그래서 미분하면 와이햇 빼기 y로 깔끔합니다. 단, 숨은 층의 ReLU 같은 activation은 분포와 상관없습니다. 마지막 층만입니다. 다음 쪽에서 네 분포의 짝을 표로 봅니다.}
\end{frame}

% --- 짝 맞추기
\begin{frame}{출력 함수와 손실은 짝으로 쓴다}
\vfill
\centering
\renewcommand{\arraystretch}{1.6}
{\large\begin{tabular}{l|c|l|l}
$y$ & 출력 $g(z)$ & 손실 & PyTorch\\ \hline
실수 & 그대로 & 오차 제곱 & \texttt{MSELoss}\\
0 / 1 & $\sigma$ & 이진 교차 엔트로피 & \texttt{BCEWithLogitsLoss}\\
$K$ 개 중 하나 & softmax & 교차 엔트로피 & \texttt{CrossEntropyLoss}\\
개수 & $e^{z}$ & 포아송 & \texttt{PoissonNLLLoss}
\end{tabular}}\\[1.0em]
{\large 짝이 맞으면 기울기 $= \hat y - y$ \quad \kq{짝이 어긋나면} ($\sigma$ + 오차 제곱): $(\hat y - y)\,\hat y(1-\hat y)$}\\[0.3em]
{\small PyTorch 의 \texttt{WithLogits} · \texttt{CrossEntropyLoss} 는 $g$ 를 손실 안에서 씌운다 $\to$ 마지막 층은 $z$ 만 낸다}
\vfill
\note{그래서 신경망을 만들 때 출력 함수와 손실은 짝으로 고릅니다. y가 실수면 출력은 그대로, 손실은 오차 제곱, PyTorch의 MSELoss입니다. 0과 1이면 시그모이드와 이진 교차 엔트로피, BCEWithLogitsLoss입니다. 여러 클래스면 소프트맥스와 교차 엔트로피, CrossEntropyLoss입니다. 개수면 e의 z와 포아송 손실, PoissonNLLLoss입니다. 짝이 맞으면 기울기가 와이햇 빼기 y로 깔끔합니다. 지난 시간에 본 것처럼 시그모이드에 오차 제곱을 쓰면 시그모이드의 미분 와이햇 곱하기 1 빼기 와이햇이 더 곱해져서, 크게 틀렸는데도 와이햇이 0이나 1 근처면 기울기가 거의 0이 되어 학습이 멈춥니다. 참고로 PyTorch의 WithLogits와 CrossEntropyLoss는 시그모이드나 소프트맥스를 손실 안에서 씌우기 때문에, 신경망의 마지막 층은 z만 내보내면 됩니다.}
\end{frame}

% --- 짝이 어긋나면: 오차 제곱 vs 교차 엔트로피
\begin{frame}{짝이 어긋나면: $\sigma$ 에 오차 제곱 vs 교차 엔트로피}
\vfill
\centering
{\large $z = w \cdot x$, \quad $\hat y = \sigma(z)$, \quad $d\hat y/dz = \hat y(1-\hat y)$}\\[0.2em]
\begin{columns}[T]
\begin{column}{0.47\textwidth}
\centering
{\large\kq{오차 제곱} (Squared Error)}\\[0.3em]
$\displaystyle L = \tfrac12(\hat y - y)^2$\\[0.3em]
$\displaystyle \frac{dL}{d\hat y} = \hat y - y$\\[0.3em]
$\displaystyle \frac{dL}{dw} = (\hat y - y)\cdot\textcolor{ksared}{\hat y(1-\hat y)}\cdot x$\\[0.3em]
{\small $\hat y(1-\hat y)$ 가 \kq{남는다}}
\end{column}
\begin{column}{0.47\textwidth}
\centering
{\large\kb{교차 엔트로피} (Cross Entropy)}\\[0.3em]
$\displaystyle L = -\big[\,y\log\hat y + (1-y)\log(1-\hat y)\,\big]$\\[0.3em]
$\displaystyle \frac{dL}{d\hat y} = \frac{\hat y - y}{\textcolor{ksared}{\hat y(1-\hat y)}}$\\[0.3em]
$\displaystyle \frac{dL}{dw} = (\hat y - y)\cdot x$\\[0.3em]
{\small $\hat y(1-\hat y)$ 가 \kb{지워진다}}
\end{column}
\end{columns}\vskip0.4em
\renewcommand{\arraystretch}{1.15}
{\large\begin{tabular}{l|c}
\multicolumn{2}{l}{정답 $y = 1$ 인데 $\hat y = 0.01$ (크게 틀림) 일 때 기울기 $(dL/dw) / x$}\\ \hline
오차 제곱 & $(0.01 - 1) \times 0.01 \times 0.99 = -0.0098$\\
교차 엔트로피 & $0.01 - 1 = -0.99$ \quad ($100$ 배)
\end{tabular}}
\vfill
\note{짝이 어긋나면 어떻게 되는지 식으로 직접 봅시다. 출력은 시그모이드, z는 w 점곱 x이고, 시그모이드의 미분은 와이햇 곱하기 1 빼기 와이햇입니다. 왼쪽은 오차 제곱입니다. 2분의 1 와이햇 빼기 y의 제곱을 와이햇으로 미분하면 와이햇 빼기 y입니다. 연쇄 법칙으로 시그모이드 미분과 x를 곱하면, 와이햇 빼기 y 곱하기 와이햇 곱하기 1 빼기 와이햇 곱하기 x입니다. 시그모이드 미분이 그대로 남습니다. 오른쪽은 교차 엔트로피, 베르누이의 마이너스 log p입니다. 와이햇으로 미분하면 마이너스 y 나누기 와이햇 더하기 1 빼기 y 나누기 1 빼기 와이햇이고, 통분하면 와이햇 빼기 y 나누기 와이햇 곱하기 1 빼기 와이햇입니다. 분모에 시그모이드 미분이 생겨서, 연쇄 법칙으로 곱하면 정확히 지워지고 와이햇 빼기 y 곱하기 x만 남습니다. 숫자로 보면 차이가 큽니다. 정답이 1인데 와이햇이 0.01로 크게 틀렸을 때, 오차 제곱의 기울기는 마이너스 0.0098, 교차 엔트로피는 마이너스 0.99로 100배 차이입니다. 오차 제곱은 크게 틀렸는데도 거의 움직이지 않습니다. 그래서 출력 함수와 손실을 짝으로 씁니다.}
\end{frame}

\section{연습}
% --- 문제풀이
\begin{frame}{문제풀이}
\vfill
\begin{enumerate}
\item 베르누이에서 $\phi = 0.8$ 이면 $z$ 는? \; ($\log 4 \approx 1.39$)\vskip0.6em
\item 보험 모델 $\hat y = e^{\,-4.137 + 0.0260\,x}$ ($x$ = 등급). 등급이 20점 오르면 사고 횟수는 몇 배? \; ($e^{0.52} \approx 1.68$)\vskip0.6em
\item 클래스 2개 소프트맥스에서 $z_1 = 3$, $z_0 = 1$ 이면 $P(y = 1)$ 은? \; ($\sigma(2) \approx 0.88$)\vskip0.6em
\item 포아송 회귀에서 정답 $y = 5$, 예측 $\hat y = 3$, 입력 $x = 2$. 이 데이터 하나의 $\partial\,\text{loss} / \partial w_1$ 은?\vskip0.6em
\item 일주일 동안 가게에 온 손님 수를 예측하는 신경망. 마지막 층의 출력 함수와 손실은?
\end{enumerate}
\vfill
\note{1번 답. z는 로그 오즈, log 0.8 나누기 0.2, log 4, 약 1.39입니다. 2번 답. 20점 오르면 z가 0.52 늘고, 사고 횟수는 e의 0.52, 약 1.68배가 됩니다. 10점마다 1.30배를 두 번 곱한 것과 같습니다. 3번 답. 시그모이드 z1 빼기 z0, 시그모이드 2, 약 0.88입니다. 4번 답. 와이햇 빼기 y는 3 빼기 5, 마이너스 2, 여기에 x 2를 곱해서 마이너스 4입니다. 음수니까 w1을 키워서 예측을 늘리는 쪽으로 갑니다. 포아송이어도 로지스틱회귀와 같은 계산입니다. 5번 답. 손님 수는 개수이니 포아송, 출력 함수는 e의 z, 손실은 포아송 손실, PyTorch의 PoissonNLLLoss입니다.}
\end{frame}

\section{정리}
% --- 정리
\begin{frame}{정리}
\vfill
\centering\large\begin{tabular}{l}
\kb{핵심} \quad 모델 고르기 $=$ $y$ 의 \kb{분포} 가정 $\to$ 씌울 함수 · 손실이 따라 나온다\\[0.5em]
GLM \quad 분포 하나 고르기 + 예측 $=$ 평균 + 매개변수 $=$ 직선 $z$\\[0.5em]
가족 \quad 정규 $\to z$ \;·\; 베르누이 $\to \sigma(z)$ \;·\; 포아송 $\to e^{z}$ \;·\; 범주형 $\to$ softmax\\[0.5em]
학습 \quad 손실 $= -\log p$, 기울기는 언제나 $(\hat y - y)\,x$ \;(이유: 지수족, 보충 덱)\\[0.5em]
신경망 \quad 마지막 층 $=$ GLM, 역전파는 $\hat y - y$ 에서 시작\\[0.5em]
보험 169,349계약 \quad 등급 10점 오를 때마다 사고 $\times 1.30$, 직선과 달리 휘고 음수 없음
\end{tabular}
\vfill
\note{정리합니다. 오늘의 핵심은 모델을 고른다는 것이 y가 어떻게 생겨나는지, 어떤 분포를 따르는지 가정하는 것이고, 그 가정 하나에서 직선에 씌울 함수와 손실이 따라 나온다는 것입니다. GLM은 y의 분포 하나를 고르고, 예측을 그 분포의 평균으로, 분포의 매개변수를 직선 z로 두는 요리법입니다. 분포에 따라 예측은 z 그대로, 시그모이드, e의 z, 소프트맥스가 되고, 각각 선형회귀, 로지스틱회귀, 포아송 회귀, 소프트맥스 회귀입니다. 손실은 마이너스 log p이고, 기울기는 언제나 와이햇 빼기 y 곱하기 x입니다. 왜 그런지는 보충 덱의 지수족에서 볼 수 있습니다. 신경망의 마지막 층도 GLM이고, 역전파는 와이햇 빼기 y에서 시작합니다.}
\end{frame}

% --- 참고 문헌
\begin{frame}{참고 문헌}
\vfill
\footnotesize
\begin{itemize}
\item Ng, A., \& Ma, T. CS229 Lecture Notes, Chapter 3: Generalized Linear Models. Stanford. \paperlink{https://cs229.stanford.edu/main_notes.pdf}
\item Ng, A. (2018). Lecture 4: Perceptron \& Generalized Linear Model, CS229 Autumn 2018. \paperlink{https://www.youtube.com/watch?v=iZTeva0WSTQ}
\item McCullagh, P., \& Nelder, J. A. (1989). \textit{Generalized Linear Models} (2nd ed.). Chapman \& Hall.
\item Dutang, C., \& Charpentier, A. CASdatasets: freMTPL2freq (프랑스 자동차 책임보험 678,013계약), OpenML 41214. \paperlink{https://www.openml.org/d/41214}
\item Bortkiewicz, L. (1898). \textit{Das Gesetz der kleinen Zahlen}. \quad Clarke, R. D. (1946). An application of the Poisson distribution. \textit{Journal of the Institute of Actuaries}, 72(3).
\item Shaw, L. (2025). V1 and V2 impact sites in London (CC BY 4.0), Zenodo. \paperlink{https://zenodo.org/records/14883986} \quad Shaw, L. P., \& Shaw, L. F. (2019). The flying bomb and the actuary. \textit{Significance}, 16(5).
\item 지도: \copyright\ OpenStreetMap contributors (ODbL). \paperlink{https://www.openstreetmap.org/copyright}
\item Fisher, R. A. (1936). The use of multiple measurements in taxonomic problems. \textit{Annals of Eugenics}, 7(2).
\item lichess.org open database, 2013년 1월 (CC0). \paperlink{https://database.lichess.org}
\end{itemize}
\vfill
\note{오늘 내용은 스탠퍼드 CS229 강의 노트의 3부, 일반화 선형 모델을 바탕으로 했습니다. 앤드류 응의 2018년 강의 4강 영상에서 같은 내용을 볼 수 있습니다. 보험 데이터는 프랑스 자동차 보험 공개 데이터 freMTPL2이고, 기병대와 폭탄은 보르트키에비치와 클라크의 원래 기록입니다. 붓꽃은 피셔의 1936년 데이터, 체스는 lichess 공개 데이터베이스입니다.}
\end{frame}
\end{document}
